problem:  
Let \( (X,d,\mathfrak{m}) \) be a **metric measure space** satisfying the **measure contraction property \( \mathrm{MCP}(K,N) \)** for some real \( K \) and finite \( N>1 \), but not necessarily a \( \mathrm{CD}(K,N) \) or \( \mathrm{RCD}(K,N) \) space. It is known that certain Carnot groups, such as the Heisenberg group with its sub-Riemannian distance and Haar measure, satisfy an \( \mathrm{MCP}(K,N) \) condition for some \( (K,N) \)—even though they fail all finite-dimensional \( \mathrm{CD}(K,N) \) conditions.  

**Question:**  
Can there exist an \( \mathrm{MCP}(K,N) \) metric measure space such that, at \( \mathfrak{m} \)-almost every point, all pointed measured Gromov–Hausdorff tangent cones are isometric to a fixed *nonabelian* Carnot group endowed with its Carnot–Carathéodory metric and Haar measure?  

Equivalently, can there exist a globally nonhomogeneous metric measure space that still has an *almost everywhere homogeneous sub-Riemannian infinitesimal structure* compatible with an \( \mathrm{MCP}(K,N) \) lower curvature bound, or do existing rigidity mechanisms force any such structure to be globally homogeneous (i.e., the Carnot group itself)?  

Why is it a "good" problem:  
This problem probes the boundary of our understanding between synthetic curvature conditions and the infinitesimal geometry of sub-Riemannian spaces. It consciously shifts from the too-rigid \( \mathrm{CD}(K,N) \) framework to the weaker \( \mathrm{MCP}(K,N) \) where Carnot groups can appear, asking whether one can construct or obstruct **nontrivial, inhomogeneous spaces** with **almost everywhere Carnot-type tangents**. Answering this would advance both measure-theoretic and geometric rigidity theory: a positive result would reveal a new class of synthetic geometries whose infinitesimal behavior is entirely sub-Riemannian, while a negative result would extend the known rigidity phenomena in curvature-dimension theory beyond the Riemannian world. The problem combines the simplicity of a local-to-global question with profound implications for how curvature bounds, homogeneity, and noncommutative geometric structures interact.